While language model (LM)-powered chatbots and generative search engines excel at answering concrete queries, discovering information in the terrain of \textit{unknown unknowns} remains challenging for users. %Such serendipitous discovery, however, is crucial for learning and complex information-seeking scenarios. 
To emulate the common educational scenario where children/students learn by listening to and participating in conversations with their parents/teachers, we create \textbf{Collaborative STORM} (\textbf{\system}). \footnote{Our resources and code are released at \url{https://github.com/stanford-oval/storm}.}
Unlike QA systems that require users to ask all the questions, \system lets users observe and occasionally steer the discourse among several LM agents. The agents ask questions on the user's behalf, allowing the user to discover \textit{unknown unknowns} serendipitously. To facilitate user interaction, \system assists users in tracking the discourse by organizing the uncovered information into a dynamic mind map, ultimately generating a comprehensive report as takeaways. For automatic evaluation, we construct the \data dataset by collecting real information-seeking records with user goals. \system outperforms baseline methods on both discourse trace and report quality. In a further human evaluation \footnote{A live research preview is now available publicly at \url{https://storm.genie.stanford.edu}}, 70\% of participants prefer \system over a search engine, and 78\% favor it over a RAG (Retrieval Augmented Generation) chatbot.

%, an information-seeking assistance system that emulates collaborative discourse among the user and LM agents with different roles. To orchestrate multiple agents, \system proposes a novel collaborative discourse protocol with different utterance intents and a turn management strategy that encourages the \moderator role to raise new questions when the discourse becomes repetitive or away from the user's goal. 